\section{Conclusion}
\label{sec:conclusion}

This paper challenges a foundational assumption in RAG engineering: that
improving retrieval quality will improve answer faithfulness. Across 6,050+
queries spanning two datasets, two model families, and four retrieval
configurations, we demonstrate that this assumption is unreliable and can be
inverted. Naive post-retrieval refinement improves local matching scores
while degrading faithfulness by up to 15.9 percentage points; and
domain-mismatched retrieval can make RAG less faithful than zero-shot
prompting.

The unifying explanation is context coherence. Faithfulness depends on
whether the retrieved passages form a mutually supportive evidence set, not
only on whether each passage is individually relevant. HCPC-v2 helps because
it respects this structure: it intervenes conservatively through AND-gated
detection, protects high-ranked passages, and restores contiguity when
refinement occurs. CCS provides a computationally cheap retrieval-time
signal for whether the evidence set is coherent enough to support faithful
generation.

Three findings have immediate practical value. First, chunk size is the
dominant design parameter for faithfulness among the factors tested
(90.1\% of between-factor variance on SQuAD),
outweighing prompt engineering by a factor of 450$\times$. Second,
post-retrieval refinement should be applied selectively based on context
coherence, not as a default post-processing step. Third, coherence should be
validated before deploying RAG in any specialized domain where the encoder is
not matched to the corpus.

\paragraph{Future directions.}
Several extensions would strengthen the generality of these findings.
Testing at larger model scales (70B+) would reveal whether the refinement
paradox persists when generators have more parametric capacity.
Domain-specific embedding models (e.g., BioLinkBERT, PubMedBERT) would test
whether fixing the encoder-corpus alignment eliminates the PubMedQA failure
mode. Multi-hop and conversational QA settings would test whether coherence
plays an even larger role when the generator must synthesize evidence across
reasoning steps. Finally, extending CCS from a diagnostic metric to an
active retrieval objective---optimizing the retrieved set for coherence
rather than individual relevance---could yield a new class of
coherence-aware retrievers.
